% Native marginal f/x notes: https://github.com/sillsdev/ptx2pdf/issues/1258
% Four verses, eight notes, inner side: the smallest fixture where the margin
% note solver left a sub-sp residual and a note's y alternated by 1sp between
% passes, so PTXprint always hit its five-pass limit. The solver is in Python,
% so plain xetex cannot show this; `make converge` runs RunJob.run_xetex() on
% this file three times and fails unless each run settles within three passes
% (two after the fix). `make reproduce` checks for the old jitter instead.
% Both need ptxprint importable, e.g. make converge PYTHON=.venv/bin/python
\def\pluginlist{}
\input paratext2.tex
\stylesheet{usfm.sty}
\stylesheet{ptx2pdf.sty}
\stylesheet{test.sty}
\def\regular{"[Charis-Regular.ttf]"}
\def\bold{"[Charis-Bold.ttf]"}
\def\italic{"[Charis-Italic.ttf]"}
\def\bolditalic{"[Charis-BoldItalic.ttf]"}
\PaperHeight=210mm
\PaperWidth=148mm
\MarginUnit=15mm
\def\TopMarginFactor{1}
\def\BottomMarginFactor{1}
\def\SideMarginFactor{2.2}
\BodyColumns=1
\FontSizeUnit=1pt
\def\LineSpacingFactor{1}
\def\VerticalSpaceFactor{1}
\hyphenpenalty=10000
\begingroup
\catcode`\@=11
\marginaln@te{f}
\marginaln@te{x}
\global\marginnotesgap=3mm
\global\marginnoteswidth=27mm
\gdef\MarginNoteSide{inner}
\endgroup
\nolotwritetrue
\ptxfile{test.usfm}
\bye
